\section{Introduction}
\label{sec:introduction}

Language models increasingly act through external tools, modify persistent
state, and operate over multiple turns.  Systems such as ReAct and ToolLLM
established that models can interleave reasoning with actions and learn broad
tool interfaces \citep{yao2023react,qin2024toolllm}.  Stateful benchmarks and
reinforcement-learning frameworks now evaluate complete goals rather than
isolated function selection
\citep{yao2024taubench,barres2025tau2bench,jiang2025verltool,chai2025rlfactory,xi2025agentgymrl}.
Yet tool execution in deployment is not stable: authentication expires,
timeouts obscure whether writes committed, responses become stale or
truncated, and schemas drift.  These events change observations and state
transitions during an episode, so recovery requires more than retrying a
failed call.

Recent methods respond with noise curricula, failure-conditioned rollout
generation, and explicit budget allocation
\citep{chen2026noisyagent,zhang2026fissiongrpo,hu2026infotree}.  Such systems
are necessarily composite.  Failure branching changes the number and
distribution of rollouts; recovery branches may use a different advantage
group; extra turns change generated-token and tool-call budgets; and
worst-group weighting changes the optimization target.  Consequently, an
end-to-end improvement does not by itself establish that branching or
minimax weighting is the active mechanism.  This is particularly important
for tool-agent RL, where small frozen evaluation sets and floor or ceiling
effects can make a single seed look decisive.

We study this identification problem through \method{}, a controlled recovery
training system with typed failure-prefix branching, executable progress
rewards, a bounded recovery budget, and fault-group minimax weighting.  Our
question is not whether this collection can be made to produce an isolated
positive run.  It is whether its proposed mechanisms retain a reproducible
behavioral advantage once the rollout estimator and interaction budget are
matched.

The evaluation is frozen before the formal comparison.  We first calibrate a
base adapter into a discriminative acceptance band, then freeze task--fault
manifests, executable state verifiers, arm definitions, and stopping rules.
Every arm is evaluated on the same episodes, enabling task-paired inference.
The formal round uses Qwen3.5-9B with LoRA and exactly the preregistered seeds
23, 29, and 43.  Three directional hypotheses isolate \method{} versus GRPO,
no-minimax, and no-fission.  A hypothesis requires the same sign at all three
seeds and a pooled paired-bootstrap confidence interval excluding zero; a
separate regression gate can veto the round.

The result is informative precisely because it is negative.  On the primary
32-episode suite, \method{} scores 16, 15, and 17 successes across the three
seeds.  Its paired differences versus GRPO are $-1,+1,0$; versus no-fission,
$0,-3,+1$; and versus no-minimax, $-5,-6,-3$.  Thus no hypothesized advantage
replicates.  Moreover, the pooled interval for \method{} minus no-minimax lies
below zero, identifying minimax weighting as harmful in this regime.  The
regression gate passes, which rules out a catastrophic-run interpretation.

This work makes four contributions:

\begin{itemize}[leftmargin=*]
    \item We provide a pre-calibrated, frozen evaluation protocol for
    executable tool-recovery RL, with task-paired outcomes, explicit
    regression gates, and directional multi-seed decision rules.
    \item We construct estimator-matched ablations that equalize primary
    rollout count and advantage grouping, with a fixed recovery budget for
    fission-bearing arms, separating mechanism claims from extra-sampling
    effects.
    \item We report a complete three-seed negative result: neither the full
    \method{} arm nor failure-prefix fission has a reproducible primary-suite
    advantage, while minimax weighting is consistently detrimental.
    \item We document the boundary of the evidence.  The result concerns a
    controlled Qwen3.5-9B LoRA regime and is not an official reproduction of
    Fission-GRPO.  Planned zero-touch external evaluation could not be run
    under a clearly applicable data and software license, so no external
    transfer claim is made.
\end{itemize}

The contribution is therefore not a new state of the art.  It is evidence
about identifiability: once estimators and budgets are aligned, intuitively
appealing recovery components may have no measurable benefit, and robust
worst-group optimization may even reduce task success.
